\clearpage
\fancyhead[L]{\small HARDING. HIMMELSATLAS.}
\SourceLeaf{515}
\begin{center}\small Göttingische gelehrte Anzeigen. Stück 61. S. 601..607. 17. April 1809.\\[3pt]\rule{.72\linewidth}{.3pt}\end{center}

\textsc{K. L. Harding}. \emph{Neuer Himmelsatlas}. \emph{Erste Lieferung} mit den Blättern I. II. V und IX. Beim Verfasser und in Kommission bei \textsc{Friedrich Perthes} in Hamburg, 1809.

Mit Vergnügen zeigen wir die ersten Früchte eines Unternehmens an, auf dessen Ausführung Kenner und Liebhaber der Astronomie seit mehreren Jahren mit Verlangen gewartet haben. Die schwierigen Zeitumstände haben ihm zwar schwere Hindernisse in den Weg gelegt, seinen Fortgang aber glücklicherweise nicht aufgehalten. Das Bedürfnis nach möglichst genauen und sehr detaillierten Himmelskarten war schon früher oft zu spüren, besonders wenn Kometen erschienen. Seit der Entdeckung der Ceres und der drei anderen verschwisterten neuen Planeten ist dieses Bedürfnis für alle, die diese so bemerkenswerten Himmelskörper beobachten oder auch nur sehen wollen, so dringend geworden, dass jeder Freund der Sternkunde es nur höchst willkommen heißen kann, wenn diese Lücke geschlossen wird, und zwar so, wie es diese ersten vier Blätter zeigen. Unser verdienter Kollege hatte sich schon vor mehr als fünf Jahren, bald nach der Entdeckung der Ceres und der Pallas, zu diesem Unternehmen entschlossen. Anfangs war der Plan freilich etwas enger gefasst. Es sollten nur diejenigen Zonen des Himmels bearbeitet werden, in denen uns diese beiden Planeten erscheinen können. Die Abgrenzung dieser Zonen gab damals Anlass zu einem interessanten analytischen Problem, das Herr Professor \textsc{Gauss} in der Monatlichen Correspondenz (August 1804 [S. 106 d. B.]) gelöst hat. Herr Professor \textsc{Harding} begann seine Arbeit unter den glücklichsten Vorzeichen. Die Entdeckung eines dritten neuen Planeten ({\LatitudeSymbols οὐ τύχης ἀλλ' ἀρετῆς εὐτυχούσης ἔργον}, wie \textsc{Delambre} mit Plutarchs Worten sagt) belohnte gleichsam im Voraus die umfangreiche Arbeit, die er auf sich genommen hatte. Aber gerade diese Entdeckung und die im Jahr 1807 Herrn Dr. \textsc{Olbers} geglückte Auffindung eines vierten Planeten machten die Erweiterung des Plans schon notwendig. Die Überzeugung von der Wichtigkeit der Sache und die vielfach von anderen Astronomen geäußerten Wünsche bewogen Herrn Professor \textsc{Harding}, die ursprünglichen Beschränkungen ganz fallen zu lassen und das Unternehmen auf den gesamten Fixsternhimmel auszudehnen, soweit er für europäische Beobachter von Interesse ist und soweit angemessene Hilfsmittel zu seiner Bearbeitung vorhanden sind. Vom ersten, enger gefassten Plan sind daher nur in der Nummerierung und Anordnung der Karten Spuren geblieben. Wer sich für diese neue Darstellung des Sternhimmels interessiert, dem muss es natürlich auch am gelegensten sein, dass die Blätter mit Teilen des Zodiakus der neuen Planeten zuerst erscheinen, um dem dringendsten Bedürfnis so früh wie möglich abzuhelfen.

Herr Professor \textsc{Harding} hat \emph{alle} Sterne in seine Karten aufgenommen, für die \emph{zuverlässige} Bestimmungen oder Beobachtungen vorliegen. Außerdem hat er nach eigenen Beobachtungen eine beträchtliche Zahl von Sternen eingetragen, die noch in keinem Verzeichnis vorkommen. Diese Bestimmungen wird er dem Atlas demnächst beifügen. Bedenkt man, wie viel besonders in letzter Zeit für die bessere Kenntnis des Fixsternhimmels getan worden ist, kann man selbst beurteilen, wie sehr der neue Atlas mit der Fülle seines Inhalts alle seine Vorgänger übertreffen muss. So hat etwa das Sternbild des Wassermanns auf den Himmelskarten von \textsc{Flamsteed} 108 Sterne und auf den großen Himmelskarten von \textsc{Bode} 343. Auf den Karten von Herrn Professor \textsc{Harding}, Blatt IX und VIII, zählen wir dagegen 1444 Sterne. Blatt VIII ist zwar erst für die zweite Lieferung bestimmt, liegt uns aber schon als Stich vor. Es ist schwer zu entscheiden, ob diese Fülle selbst oder die beispiellose Gewissenhaftigkeit, mit der sie genutzt wurde, als der größe\SourceLeaf{516}re Vorzug dieser Karten gelten soll. Alle bisherigen Bearbeiter von Himmelskarten haben die Sterne stets so eingetragen, wie die für die besten gehaltenen Quellen sie angeben. Auch wenn sie, wie Herr Professor \textsc{Bode} bei seinem Atlas, durch sorgfältige Vergleiche, Vermutungen und scharfsinnige Kritik viele Fehler aller Art erfolgreich beseitigten, die bei einer so ungeheuren Menge an Zahlen in den Fixsternverzeichnissen unvermeidlich waren, mussten doch sehr viele Fehler unbemerkt bleiben und in die Karten übergehen, weil diese Mittel zu ihrer Entdeckung nicht ausreichten. Daher stehen auf den bisherigen Himmelskarten so viele Sterne, die am Himmel gar nicht stehen und nie gestanden haben. So hat Herr Professor \textsc{Harding} unter den 343 Sternen des Wassermanns bei \textsc{Bode} nicht weniger als 15 solcher Sterne bemerkt, in manchen anderen Sternbildern verhältnismäßig noch weit mehr. Das einzige Mittel, diesen Fehlern zu entgehen, bestand darin, \emph{sämtliche} Sterne vor ihrer Aufnahme in die Karten sorgfältig \emph{am Himmel selbst} zu überprüfen. Bei einem Unternehmen dieses Umfangs schien das freilich die Kräfte eines einzelnen Menschen beinahe zu übersteigen. Dennoch hielt Herr Professor \textsc{Harding} es bei der Bearbeitung eines neuen Atlas für unerlässlich. In die vorliegenden Himmelskarten wurde daher kein einziger Stern aufgenommen, von dessen wirklichem Vorhandensein an seinem Platz sich Herr Professor \textsc{Harding} nicht vorher selbst überzeugt hätte.

Aus diesen allgemeinen Bemerkungen über den Geist, in dem das Unternehmen ausgeführt wird, lässt sich leicht auf seinen wissenschaftlichen Wert schließen. Wir wollen noch einige Worte über das Äußere der Karten hinzufügen. Herr Professor \textsc{Harding} meinte, er solle ein wesentlich kleineres Format wählen als das der Karten von \textsc{Bode}, teils weil es bequemer zu gebrauchen ist, teils um den Preis nicht zu sehr erhöhen zu müssen --- beides Gesichtspunkte, die für Astronomen sehr ins Gewicht fallen. Die Blätter sind $16\frac{3}{4}$ Pariser Zoll hoch und $20\frac{1}{2}$ Zoll breit. Sie haben also nur etwa ein Drittel der Fläche der Karten von \textsc{Bode}. Bei den ersten vier Blättern, auf denen die Deklinationen nicht über 30 Grad nach Norden und Süden hinausgehen, schien es mit Recht überflüssig, eine eigentliche Projektion zu wählen. Die Grade der Rektaszension und der Deklination sind daher durchgehend gleich groß angesetzt (etwa $5\frac{6}{7}$ Pariser Linien). Das bietet bei der Benutzung noch einige Vorteile, die die geringe, selbst bei den größten Deklinationen kaum merkliche Verzerrung bei Weitem aufwiegen. Jedes Blatt umfasst so 42 Grad in Rektaszension und $34\frac{1}{2}$ Grad in Deklination. In Rektaszension reicht das erste Blatt von $359^\circ$ bis $41^\circ$, das zweite von $39^\circ$ bis $81$ usw., sodass sich an das neunte Blatt von $319^\circ$ bis $1^\circ$ wieder das erste anschließt. Das Gradnetz für Rektaszension und Deklination ist zur bequemeren Benutzung stehen geblieben. Außer diesem Netz und den Grenzen und Namen der Sternbilder enthalten die Karten nur das, was jemand auf ihnen sucht, der sie zum wirklichen Gebrauch und nicht als Sammlung hübscher Bilder anschafft. Das sind die Sterne und ihre allgemein gebräuchlichen Bezeichnungen, also die Buchstaben von \textsc{Bayer} oder die Zahlen von \textsc{Flamsteed}. Die Figuren der Sternbilder selbst sind beim heutigen Stand der Wissenschaft ganz außer Gebrauch. In den Augen eines Kenners hätten sie den Hintergrund der Karten nur verunstaltet. Man hat sie mit Recht ganz weggelassen und dadurch erst die deutliche und saubere Darstellung einer solchen Fülle von Sternen ermöglicht. Im Durchschnitt wird man auf jedem Blatt über 2000 Sterne zählen können. Wenn das Ganze also nach dem Plan des Verfassers aus 24 Blättern bestehen wird, von denen der größere Teil bereits vollständig gezeichnet ist, wird man in diesem Himmelsatlas ungefähr 50000 Sterne besitzen.

Mancher würde sich vielleicht wünschen, dass der Verfasser außer dem Netz für Rektaszension und Deklination noch ein Netz für die Längen- und Breitengrade hinzugefügt hätte, \SourceLeaf{517}vor allem um Planetenpositionen bequemer eintragen zu können, die in den astronomischen Ephemeriden gewöhnlich nur nach Länge und Breite angegeben werden. Dieses Netz wurde jedoch mit Bedacht weggelassen, damit das Gerüst schwarzer Linien den Eindruck, den die Anordnung der Sterne auf den Karten auf das Auge macht, möglichst wenig stört. Als Ersatz wird Herr Professor \textsc{Harding} in dem am Ende beizufügenden Text eine sehr handliche, sorgfältig berechnete Tabelle mitteilen. Mit ihrer Hilfe kann jeder Liebhaber dieses Netz ganz leicht mit der Feder selbst in sein Exemplar eintragen, wenn er möchte etwa in einer anderen Farbe, um beide Zwecke zugleich zu erreichen. Vielleicht wäre der Verfasser bereit, diese nützliche Tafel schon früher an geeigneter Stelle für die Besitzer der ersten Lieferung mitzuteilen.

Uns bleibt nur noch, die Teile des Himmels zu nennen, die auf den vier Blättern der ersten Lieferung dargestellt sind. Die Karten sind nach den Graden der Rektaszension und Deklination statt nach den Sternbildern eingeteilt, um möglichst viel Platz zu sparen und unnötige Wiederholungen zu vermeiden. Die meisten Sternbilder sind daher aufgeteilt. Dem kundigen Liebhaber braucht man nicht erst zu sagen, dass dies ganz gleichgültig ist. Das Blatt Nr. I. erstreckt sich in Deklination von $13^\circ$ nördlich bis $21\frac{1}{2}^\circ$ südlich; es enthält also Teile der Fische und des Widders, den größten Teil des Walfischs und einen kleinen Teil des Eridanus. Auf Blatt Nr. II. von $28\frac{1}{2}^\circ$ nördlich bis $6^\circ$ südlich findet man größtenteils den Rest des Widders und des Walfischs, fast den ganzen Stier und die Georgsharfe sowie Teile des Eridanus und des Orion. Blatt Nr. V. von $27^\circ$ nördlich bis $7\frac{1}{4}^\circ$ südlich zeigt einen großen Teil des Löwen, Teile des Kleinen Löwen und des Bechers sowie den größeren Teil der Jungfrau und des Haars der Berenike. Schließlich zeigt Blatt Nr. IV. von $5^\circ$ nördlich bis $29\frac{1}{4}^\circ$ südlich einen Teil des Steinbocks, des Luftballons, des Südlichen Fisches, der Bildhauerwerkstatt, des Walfischs, der Fische, des Pegasus und des Füllens sowie den größten Teil des Wassermanns. --- Wir schließen diese Anzeige mit dem Wunsch, dass alle, denen das Gedeihen der Wissenschaften und die Ermunterung echt deutscher Beharrlichkeit bei wahrhaft nützlichen Arbeiten am Herzen liegen, nun auch dieses Unternehmen unterstützen und aufrechterhalten, das man gar nicht genug empfehlen kann. Möge es als beste Belohnung für seinen würdigen Urheber künftig für die erhabene Wissenschaft reiche Früchte tragen.
